\clearpage
\cleardoublepage

\chapter{优化器与学习率调度器}

\section{梯度下降优化}

训练神经网络通过在减少损失的方向上更新参数来最小化损失函数。

\subsection{批量梯度下降}

使用整个数据集的梯度更新参数：
\begin{equation}
\theta_{t+1} = \theta_t - \alpha \nabla_\theta \mathcal{L}(\theta_t)
\end{equation}

\begin{itemize}
    \item \textbf{优点：} 收敛稳定，梯度精确
    \item \textbf{缺点：} 大数据集上慢，内存占用大
\end{itemize}

\subsection{随机梯度下降（SGD）}

使用单个样本的梯度更新参数：
\begin{equation}
\theta_{t+1} = \theta_t - \alpha \nabla_\theta \mathcal{L}(\theta_t; \mathbf{x}^{(i_t)}, \mathbf{y}^{(i_t)})
\end{equation}

\begin{itemize}
    \item \textbf{优点：} 快，能跳出局部最小值，可在线学习
    \item \textbf{缺点：} 梯度噪声大，需要较小的学习率
\end{itemize}

\subsection{小批量梯度下降}

使用 $B$ 个样本的批次更新：
\begin{equation}
\theta_{t+1} = \theta_t - \alpha \frac{1}{B}\sum_{i=1}^{B}\nabla_\theta \mathcal{L}(\theta_t; \mathbf{x}^{(i)}, \mathbf{y}^{(i)})
\end{equation}

这是深度学习中的标准方法。

\section{带动量的梯度下降}

动量通过累积速度来加速收敛：
\begin{align}
\mathbf{v}_{t+1} &= \beta \mathbf{v}_t + (1-\beta)\nabla_\theta \mathcal{L}(\theta_t) \\
\theta_{t+1} &= \theta_t - \alpha \mathbf{v}_{t+1}
\end{align}

\begin{itemize}
    \item $\beta \in [0, 1)$ 是动量系数（通常为 0.9）
    \item \textbf{减少}陡峭方向的振荡
    \item \textbf{加速}平缓方向的收敛
\end{itemize}

\begin{note}[Nesterov 加速梯度]
Nesterov 动量向前看：
\begin{align}
\mathbf{v}_{t+1} &= \beta \mathbf{v}_t - \alpha \nabla_\theta \mathcal{L}(\theta_t + \beta\mathbf{v}_t) \\
\theta_{t+1} &= \theta_t + \mathbf{v}_{t+1}
\end{align}
这提供了更好的收敛保证，通常训练更快。
\end{note}

\section{自适应学习率方法}

这些方法根据梯度历史调整每个参数的学习率。

\subsection{AdaGrad}

学习率与累积平方梯度成反比地调整：
\begin{align}
\mathbf{g}_t &= \nabla_\theta \mathcal{L}(\theta_t) \\
\mathbf{r}_t &= \mathbf{r}_{t-1} + \mathbf{g}_t \odot \mathbf{g}_t \\
\theta_{t+1} &= \theta_t - \frac{\alpha}{\sqrt{\mathbf{r}_t + \epsilon}} \odot \mathbf{g}_t
\end{align}

\begin{itemize}
    \item \textbf{适合稀疏特征}
    \item \textbf{学习率随时间递减}
    \item \textbf{在深度学习中可能过早停止}
\end{itemize}

\subsection{RMSProp}

使用平方梯度的指数加权移动平均：
\begin{align}
\mathbf{g}_t &= \nabla_\theta \mathcal{L}(\theta_t) \\
\mathbf{r}_t &= \beta \mathbf{r}_{t-1} + (1-\beta)\mathbf{g}_t \odot \mathbf{g}_t \\
\theta_{t+1} &= \theta_t - \frac{\alpha}{\sqrt{\mathbf{r}_t + \epsilon}} \odot \mathbf{g}_t
\end{align}

\begin{itemize}
    \item \textbf{解决了 AdaGrad 学习率递减的问题}
    \item \textbf{RNN 的常用默认值}
    \item $\beta = 0.9$ 是典型值
\end{itemize}

\subsection{Adam（自适应矩估计）}

Adam \cite{kingma2014adam} 结合了动量和 RMSProp：
\begin{align}
\mathbf{m}_t &= \beta_1 \mathbf{m}_{t-1} + (1-\beta_1)\mathbf{g}_t \quad \text{（一阶矩）} \\
\mathbf{v}_t &= \beta_2 \mathbf{v}_{t-1} + (1-\beta_2)\mathbf{g}_t^2 \quad \text{（二阶矩）} \\
\hat{\mathbf{m}}_t &= \frac{\mathbf{m}_t}{1-\beta_1^t} \quad \text{（偏差校正）} \\
\hat{\mathbf{v}}_t &= \frac{\mathbf{v}_t}{1-\beta_2^t} \\
\theta_{t+1} &= \theta_t - \frac{\alpha \hat{\mathbf{m}}_t}{\sqrt{\hat{\mathbf{v}}_t} + \epsilon}
\end{align}

\begin{properties}[Adam 的特性]
\begin{itemize}
    \item \textbf{大多数深度学习任务的默认优化器}
    \item \textbf{对超参数鲁棒}
    \item \textbf{默认参数：} $\alpha = 0.001, \beta_1 = 0.9, \beta_2 = 0.999, \epsilon = 10^{-8}$
    \item \textbf{早期时间步的偏差校正}
\end{itemize}
\end{properties}

\begin{code}
    \captionof{listing}{PyTorch 中的优化器}
    \label{code:optimizers}
\begin{lstlisting}[language=python]
import torch.optim as optim

# 带动量的 SGD
optimizer = optim.SGD(model.parameters(), lr=0.01, momentum=0.9)

# 带 Nesterov 动量的 SGD
optimizer = optim.SGD(model.parameters(), lr=0.01, momentum=0.9, nesterov=True)

# Adam
optimizer = optim.Adam(model.parameters(), lr=0.001, betas=(0.9, 0.999))

# AdamW（解耦权重衰减）
optimizer = optim.AdamW(model.parameters(), lr=0.001, weight_decay=0.01)

# RMSProp
optimizer = optim.RMSprop(model.parameters(), lr=0.01, alpha=0.99)
\end{lstlisting}
\end{code}

\subsection{AdamW（带权重衰减的 Adam）}

AdamW \cite{loshchilov2019decoupled} 将权重衰减与梯度更新解耦：
\begin{equation}
\theta_{t+1} = \theta_t - \alpha \frac{\mathbf{m}_t}{\sqrt{\mathbf{v}_t} + \epsilon} - \alpha \lambda \theta_t
\end{equation}

\begin{itemize}
    \item \textbf{与原始 Adam 不同的正确 L2 正则化}
    \item \textbf{在许多任务中泛化更好}
    \item \textbf{Transformer 的标准配置}
\end{itemize}

\subsection{AdaDelta}

AdaDelta \cite{zeiler2012adadelta} 通过将累积过去梯度的窗口限制为固定大小 $w$ 来解决 AdaGrad 单调递减学习率的问题，使用指数衰减平均：
\begin{align}
\mathbf{r}_t &= \rho \mathbf{r}_{t-1} + (1 - \rho)\mathbf{g}_t^2 \\
\Delta\boldsymbol{\theta}_t &= -\frac{\sqrt{\mathbf{s}_{t-1} + \epsilon}}{\sqrt{\mathbf{r}_t + \epsilon}}\mathbf{g}_t \\
\mathbf{s}_t &= \rho \mathbf{s}_{t-1} + (1 - \rho)(\Delta\boldsymbol{\theta}_t)^2
\end{align}
其中 $\rho$ 是衰减常数（通常为 0.95）。

\begin{properties}[AdaDelta 的特性]
\begin{itemize}
    \item \textbf{无需手动设置学习率：} 分子 $\sqrt{\mathbf{s}_{t-1} + \epsilon}$ 替代了 $\alpha$，使该方法基本无超参数
    \item \textbf{对不同梯度尺度鲁棒：} 使用累积更新自适应调整每个参数
    \item \textbf{现在很少使用：} 被 Adam 及其变体取代
\end{itemize}
\end{properties}

\subsection{自适应方法的统一视角}

\begin{table}[ht]
    \centering
    \caption{自适应优化器家族}
    \begin{tabular}{L{2cm} L{3cm} L{3cm} L{3.5cm}}
        \toprule
        \textbf{方法} & \textbf{一阶矩} & \textbf{二阶矩} & \textbf{LR 缩放} \\
        \midrule
        AdaGrad & 无 & $\mathbf{r}_t = \mathbf{r}_{t-1} + \mathbf{g}_t^2$ & $\alpha / \sqrt{\mathbf{r}_t}$ \\
        RMSProp & 无 & $\mathbf{r}_t = \rho\mathbf{r}_{t-1} + (1-\rho)\mathbf{g}_t^2$ & $\alpha / \sqrt{\mathbf{r}_t}$ \\
        AdaDelta & 无 & 与 RMSProp 相同 & $\sqrt{\mathbf{s}} / \sqrt{\mathbf{r}_t}$ \\
        Adam & $\mathbf{m}_t$（偏差校正） & $\mathbf{v}_t$（偏差校正） & $\alpha \hat{\mathbf{m}}_t / \sqrt{\hat{\mathbf{v}}_t}$ \\
        AdamW & 与 Adam 相同 & 与 Adam 相同 & + 解耦 $\lambda\theta_t$ \\
        \bottomrule
    \end{tabular}
\end{table}

\subsection{收敛速率分析}

\begin{theorem}[SGD 收敛速率]
对于平滑（$L$-Lipschitz 梯度）和 $\mu$-强凸函数 $f$，学习率 $\alpha = 1/L$ 的梯度下降满足：
\begin{equation}
f(\theta_t) - f(\theta^*) \leq \left(1 - \frac{\mu}{L}\right)^t \bigl(f(\theta_0) - f(\theta^*)\bigr)
\end{equation}
这是速率 $1 - \mu/L$ 的\textbf{线性收敛速率}。
\end{theorem}

\begin{properties}[不同优化器的收敛速率]
\begin{itemize}
    \item \textbf{凸 + 平滑：} SGD 达到 $O(1/T)$ 速率；加速方法达到 $O(1/T^2)$
    \item \textbf{非凸：} 所有方法收敛到 $O(1/\sqrt{T})$ 的驻点
    \item \textbf{Adam} 在非凸设置下不能保证收敛 \cite{reddi2018convergence}；\textbf{AdamW} 和 \textbf{AMSGrad} 修复了这个问题
    \item 在实践中，Adam 和 AdamW 初始收敛更快，但 SGD+动量通常能达到更好的最终泛化
\end{itemize}
\end{properties}

\section{梯度裁剪}

防止梯度爆炸。两种主要策略：

\subsection{值裁剪}

逐个裁剪每个梯度分量：
\begin{equation}
\tilde{g}_i = \text{clip}(g_i, -c, c) = \max(\min(g_i, c), -c)
\end{equation}
其中 $c$ 是裁剪值（通常为 0.5--5.0）。

\subsection{范数裁剪}

如果整个梯度向量的范数超过阈值，则按比例缩放：
\begin{equation}
\tilde{\mathbf{g}} = \begin{cases}
\mathbf{g} & \text{如果 } \|\mathbf{g}\| \leq c \\
\frac{c}{\|\mathbf{g}\|} \mathbf{g} & \text{如果 } \|\mathbf{g}\| > c
\end{cases}
\end{equation}
这保留了梯度的\textbf{方向}，同时限制其幅度。

\begin{properties}[何时使用梯度裁剪]
\begin{itemize}
    \item \textbf{RNN 和 Transformer：} 对于防止长序列中的梯度爆炸必不可少
    \item \textbf{微调预训练模型：} 防止大梯度更新造成的失稳
    \item \textbf{范数裁剪}通常优于值裁剪（保留梯度方向）
    \item 典型最大范数：1.0（Transformer）到 5.0（RNN）
    \item PyTorch 中：\texttt{torch.nn.utils.clip\_grad\_norm\_}(model.parameters(), max\_norm=1.0)
\end{itemize}
\end{properties}

\section{学习率调度器}

学习率调度器在训练过程中调整学习率。

\subsection{步进衰减}

\begin{equation}
\alpha_t = \alpha_0 \cdot \gamma^{\lfloor t / T \rfloor}
\end{equation}
其中 $\gamma$ 是衰减因子（如 0.5），$T$ 是周期。

\subsection{指数衰减}

\begin{equation}
\alpha_t = \alpha_0 \cdot e^{-\lambda t}
\end{equation}

\subsection{余弦退火}

\begin{equation}
\alpha_t = \alpha_{min} + \frac{1}{2}(\alpha_{max} - \alpha_{min})\left(1 + \cos\left(\frac{t\pi}{T}\right)\right)
\end{equation}

\begin{figure}[htbp]
    \centering
    \begin{tikzpicture}[scale=0.8]
        \begin{axis}[
            xlabel=Epoch,
            ylabel=学习率,
            width=10cm,
            height=5cm,
            legend pos=north east
        ]
            % 步进衰减
            \addplot[blue, const plot, mark=*] coordinates {
                (0, 0.01)
                (30, 0.005)
                (60, 0.0025)
                (90, 0.00125)
            };
            \addlegendentry{步进衰减}
            
            % 指数衰减
            \addplot[red, smooth] {0.01 * exp(-0.02*x)};
            \addlegendentry{指数衰减}
            
            % 余弦退火
            \addplot[green, smooth] {(0.001 + 0.5*(0.01-0.001)*(1+cos(pi*x/100))};
            \addlegendentry{余弦退火}
        \end{axis}
    \end{tikzpicture}
    \caption{学习率调度器}
    \label{fig:lr_schedules}
\end{figure}

\subsection{预热}

在开始时逐渐增加学习率：
\begin{equation}
\alpha_t = \alpha_{max} \cdot \frac{t}{T_{warmup}} \quad \text{当 } t < T_{warmup}
\end{equation}

\begin{itemize}
    \item \textbf{防止开始时的大梯度}
    \item \textbf{对 Transformer 很重要}
    \item \textbf{常见：500-10000 预热步数}
\end{itemize}

\subsection{预热 + 余弦衰减}

现代 Transformer 训练的标准调度结合了预热和余弦退火：
\begin{equation}
\alpha_t = \begin{cases}
\alpha_{max} \cdot \frac{t}{T_{warmup}} & \text{如果 } t < T_{warmup} \\
\alpha_{min} + \frac{1}{2}(\alpha_{max} - \alpha_{min})\left(1 + \cos\left(\frac{(t - T_{warmup})\pi}{T_{total} - T_{warmup}}\right)\right) & \text{否则}
\end{cases}
\end{equation}

这种预热加余弦的调度是许多 Transformer 训练方法的默认配置；见 \cite{loshchilov2016sgdr,goyal2017accurate}。

\subsection{ReduceLROnPlateau}

监控指标（通常是验证损失），当改进停滞时降低学习率：
\begin{equation}
\alpha_t = \alpha_{t-1} \cdot \gamma \quad \text{如果 } \text{连续 } p \text{ 个 epoch 无改进}
\end{equation}
其中 $\gamma$ 是衰减因子（通常为 0.1），$p$ 是耐心值（通常为 5--10 个 epoch）。

\begin{itemize}
    \item \textbf{自适应：} 无需预先定义总 epoch 数
    \item \textbf{当训练时间未知时有用}（如早停）
    \item 常用于训练 CNN 进行图像分类
\end{itemize}

\subsection{OneCycleLR}

Smith \cite{smith2018disciplined} 提出，OneCycleLR 使用单个周期，分三个阶段：
\begin{enumerate}
    \item \textbf{增加} LR 从 $\alpha_{min}$ 到 $\alpha_{max}$（约前 30\% 的步数）
    \item \textbf{减少} LR 从 $\alpha_{max}$ 到 $\alpha_{min}$（约接下来 60\% 的步数）
    \item \textbf{微调：} 进一步减少到 $\alpha_{min}/100$（最后约 10\% 的步数）
\end{enumerate}
可选地，动量遵循逆模式（LR 低时动量高）。

\begin{properties}[OneCycleLR 的优势]
\begin{itemize}
    \item \textbf{收敛更快：} 通常用更少的 epoch 达到相同精度
    \item \textbf{正则化效果：} 高到低的 LR 周期起到隐式正则化的作用
    \item \textbf{无需预热：} 内置于调度中
    \item 在 fastai 和迁移学习任务中很受欢迎
\end{itemize}
\end{properties}

\section{优化器选择}

\begin{table}[ht]
    \centering
    \caption{优化器选择指南}
    \begin{tabular}{L{2.5cm} L{3cm} L{3.5cm}}
        \toprule
        \textbf{优化器} & \textbf{适用场景} & \textbf{备注} \\
        \midrule
        SGD + 动量 & 图像分类 & 泛化最好 \\
        Adam & 通用 & 默认选择 \\
        AdamW & Transformer、NLP & 微调最佳 \\
        RMSProp & RNN & 循环模型适用 \\
        \bottomrule
    \end{tabular}
\end{table}

\section{本章小结}

\begin{enumerate}
    \item 梯度下降是神经网络优化的基础
    \item 动量通过累积速度加速收敛
    \item 自适应方法（AdaGrad、RMSProp、AdaDelta、Adam）调整每个参数的学习率
    \item Adam 结合动量和自适应学习率；AdamW 添加正确的权重衰减
    \item 梯度裁剪防止梯度爆炸（优先使用范数裁剪）
    \item 学习率调度器（步进、余弦、 plateau、OneCycleLR）改善收敛
    \item 预热对 Transformer 训练很重要
    \item SGD+动量通常泛化更好；Adam 收敛更快
\end{enumerate}

\section{扩展阅读}

\begin{itemize}
    \item \cite{zeiler2012adadelta} —— ADADELTA：一种自适应学习率方法
    \item \cite{pascanu2013difficulty} —— 训练循环神经网络的困难
    \item \cite{kingma2014adam} —— Adam：一种随机优化方法
    \item \cite{loshchilov2016sgdr} —— 带热重启的随机梯度下降
    \item \cite{reddi2018convergence} —— 关于 Adam 的收敛性及以后
    \item \cite{smith2018disciplined} —— 神经网络超参数的有原则方法
    \item \cite{loshchilov2019decoupled} —— 解耦权重衰减正则化
    \item \cite{paszke2019pytorch} —— PyTorch：一种命令式风格的高性能深度学习库
\end{itemize}

\section{习题}

\begin{enumerate}
    \item \textbf{动量更新：} 从 $\mathbf{v}_{t+1} = \beta \mathbf{v}_t + (1-\beta)\mathbf{g}_t$ 开始，解释动量如何与普通 SGD 相比平滑噪声梯度。

    \item \textbf{Adam 偏差校正：} 为什么 Adam 中需要偏差校正的矩估计 $\hat{\mathbf{m}}_t$ 和 $\hat{\mathbf{v}}_t$？如果省略它们，前几次迭代会发生什么？

    \item \textbf{AdamW vs Adam：} 展示在使用 Adam 时，解耦权重衰减与在损失上添加 L2 惩罚有何不同。为什么 AdamW 通常更适合 Transformer？

    \item \textbf{梯度裁剪：} 假设梯度向量的范数为 12，裁剪阈值为 3。写出范数裁剪后的梯度，并解释保留了什么特性。

    \item \textbf{余弦退火：} 绘制在一个训练周期中余弦学习率衰减的定性形状。为什么大的初始学习率可以帮助优化？

    \item \textbf{调度器选择：} 为以下场景选择调度器并说明理由：长期预训练运行、未知停止时间、快速迁移学习微调。
\end{enumerate}

